\chapter{Two-clock cohomology, original integral forcing, and a
source-complete finite certificate}
{\small\sffamily Source: \nolinkurl{repository/collatz_reconstruction/research_program/clocked_support_20260914/note.md}.\par}


14 September 2026. Collatz-only continuation of
\nolinkurl{KokunoYumeto/collatz-workbench}, authenticated source head
\nolinkurl{6a30f1dc23110fae0abbe3b3bd5e60072a842385}.

The contribution retains the workbench's original two-clock weight, not
a new normalization. It gives its marked cyclic complex, exact
integral-arithmetic specialization, component cokernel, completion maps,
and a clock-preserving lift of the previously constructed descent
retractions. A reversible source database closes the previously retained
403-rise sector and gives an all-length 678-rise exclusion. The actual
infinite residual source remains.

\section{0. Source relationships and notation}

\textbf{{[}W{]}}
\nolinkurl{research_companion/chapters/01_weighted_path_and_finite_residue_actions.tex},
Git blob \nolinkurl{f0142c0d4c237ed068a0762a08c6622df5dbcfee}, in the
same authenticated commit. The first-return category, odd-block
expansion, two-clock weights, endpoint series, finite-support transfer
operator, and completed inverse are already present there. This
continuation does not claim those constructions as new. The inspected
source passages are lines 1--155 and 260--555.

\textbf{{[}C{]}}
\nolinkurl{research_program/cycle_relative_cohomology_20260914/note.md},
the relative actual-orbit complex and exact descent retractions.

\textbf{{[}S{]}}
\nolinkurl{research_program/residual_splice_20260914/note.md} and
\nolinkurl{integral_blocks.md}, the complete first-descent fibers, root
splicing, original integral comparison and retained 403-rise sector. The
source note has Git blob
\nolinkurl{b31acf3b485b9565b1bbf878f1e1a98d5e704ad8}.

All paths above are relative to \nolinkurl{collatz_reconstruction/}.
Original forcing is +1 unless a control is expressly labelled otherwise.
Nothing here supplies a morphism from the entire Zeta theta quotient to
the entire Collatz orbit complex, or imports an RH conclusion. The
finite support/Hurwitz relationship already recorded in the repository
remains available. The new arrows below begin with the actual Collatz
path operator of {[}W{]}.

Put \SourceMath{X=\{1,3,5,\ldots\},\qquad T(n)=\frac{3n+1}{2^{a(n)}},
 \qquad a(n)=\nu_2(3n+1),\qquad R=\mathbb Z[u,v].}{} The weight of the
original edge at n is \SourceMath{w_n=u v^{a(n)-1}.}{C1} For a path of m
odd returns and exponent sum A it is exactly \(u^m v^{A-m}\). This is
{[}W{]}'s convention. The explicit expansion of exponent a into one
shortened odd step and a−1 shortened even steps recovers both clocks and
has the inverse obtained by grouping at odd returns.

The variable v in R is a clock variable; the symbol \(V_n\) below is an
original vertex basis vector. No original integer is rescaled in these
definitions.

\section{1. The relative complex comes from the original transfer
operator}

Let E and V be free R-modules on the original nonbase edges \(E_n\) and
vertices \(V_n\), respectively, for n in X excluding 1. Define
\SourceMath{d_{u,v} E_n=V_n-w_nV_{T(n)},\qquad V_1=0.}{C2} Give E degree
zero and V degree one, and write
\SourceMath{\mathcal K^\bullet_{u,v}=[E\xrightarrow{d_{u,v}}V].}{} At
u=v=1, this is precisely {[}C{]}'s relative chain complex with degree
reindexing \(\mathcal K^i=C_{1-i}\): the map on every original edge and
vertex is the identity. Thus H\^{}1 of this reindexed complex is the old
H\_0, and H\^{}0 is the old H\_1. This states the degree correspondence
explicitly.

The transfer formula of {[}W{]}, on finite-support vectors, is
\SourceMath{(Lf)(y)=u\sum_{T(x)=y}v^{a(x)-1}f(x).}{} Each input basis
vector has exactly one image, so \(L V_n=w_nV_{T(n)}\). Identifying
\(E_n\) with \(V_n\) gives \(d=I-L\), literally. On the original module
including 1, \(R V_1\) is invariant because \(L V_1=uvV_1\). Passing to
the specified quotient by this submodule gives (C2). On complexes, only
the subcomplex \SourceMath{[R E_1\xrightarrow{1-uv}R V_1]}{} is removed.
No other source or zero coefficient is used as a reason to remove a
basis label.

\subsection{Proposition 1. Polynomial injectivity and original path
boundaries}

The map d in (C2) is injective. For a finite actual path
\(n_0\to\cdots\to n_m\), stopping on its first hit of 1, put
\SourceMath{W_0=1,\qquad W_j=\prod_{i<j}w_{n_i},\qquad
 P_p=\sum_{j=0}^{m-1}W_j E_{n_j}.}{} Then
\SourceMath{dP_p=V_{n_0}-W_mV_{n_m}.}{C3}

\textbf{Proof.} The two terms belonging to each intermediate vertex
cancel with their identical original prefix monomial, giving (C3). For
injectivity, identify E and V by labels and write \(L=uS\), where
\(S V_n=v^{a(n)-1}V_{T(n)}\). In \((I-uS)f=0\), comparison of the
constant coefficient in u gives f\_0=0; the subsequent coefficients give
f\_j=Sf\_\{j-1\}=0. A polynomial vector has finitely many coefficients,
all zero. This works on the declared finite-support modules, without any
convergence assumption. ∎

The original graph's closed-cycle kernel at u=v=1 is therefore not
obtained by declaring the generic polynomial kernel to be nonzero. The
cycle information is retained in the polynomial cokernel and in its
explicitly computed specialization below.

\section{2. A faithful marked cyclic complex}

Keep a nonempty original labelled exponent word \(p=(a_1,\ldots,a_m)\).
Define \SourceMath{A_j=\sum_{i=1}^j a_i,\quad A_0=0,\quad
 w_i=uv^{a_i-1},\quad W_j=\prod_{i=1}^j w_i,}{}
\SourceMath{W=W_m=u^m v^{A_m-m},\qquad
 N_p(u,v)=\sum_{i=1}^m W_{i-1}.}{C4} The position-circle boundary sends
edge i to vertex i minus \(w_i\) times vertex i+1, cyclically. Its
algebraic dual on the finite free modules is
\SourceMath{(J_ph)_i=h_i-w_i h_{i+1},\qquad h_{m+1}=h_1.}{C5} The
duality is the original evaluation pairing:
\(\langle d e,h\rangle=\langle e,J_ph\rangle\). Neither (C5) nor (C2) is
identified with the integer matrix B before the specialization in
Section 3.

\subsection{Theorem 2. Explicit polynomial comparison, both homotopies,
and forcing}

The maps
\SourceMath{\phi_0(h)=h_1,\qquad \phi_1(b)=\sum_{i=1}^mW_{i-1}b_i}{}
give a chain map
\SourceMath{[R^m\xrightarrow{J_p}R^m]\longrightarrow[R\xrightarrow{1-W}R].}{C6}
A section is
\SourceMath{(\psi_0t)_1=t,\qquad (\psi_0t)_i=\left(\prod_{k=i}^m w_k\right)t\ (i\ge2),
 \qquad \psi_1t=(t,0,\ldots,0).}{} Define \SourceMath{(Hb)_1=0,\qquad
 (Hb)_i=\sum_{j=i}^m\left(\prod_{k=i}^{j-1}w_k\right)b_j\quad(i\ge2).}{C7}
They satisfy
\SourceMath{\phi_1J_p=(1-W)\phi_0,\quad J_p\psi_0=\psi_1(1-W),\quad
 \phi_0\psi_0=\phi_1\psi_1=I,}{}
\SourceMath{J_pH=I-\psi_1\phi_1,\qquad HJ_p=I-\psi_0\phi_0.}{C8}
Consequently the marked cohomology is
\SourceMath{H^0=0,\qquad H^1\cong R/(1-W),\qquad [b]\longmapsto[\phi_1(b)].}{C9}
For the unit forcing vector its mark is \([N_p]\). The exact inverse of
the forced equation is
\SourceMath{(1-W)t=\phi_1(b),\qquad h=\psi_0t+Hb.}{C10} All maps and
identities are over R, and remain valid after every specified
coefficient homomorphism.

\textbf{Proof.} Multiplication of (C5) by W\_\{i−1\} telescopes to
\((1-W)h_1\), proving the first chain equation. In J\_pψ\_0, all rows
except the first cancel; the first is (1−W)t. The section identities
follow from the first coordinate. For i≥2, the tail sum (C7) obeys
\((Hb)_i-w_i(Hb)_{i+1}=b_i\), including the final row with \((Hb)_1=0\).
The first row is \(-\sum_{i\ge2}W_{i-1}b_i=b_1-\phi_1(b)\). This proves
the first homotopy. Applying the same telescope to J\_ph in each tail
gives \(h_i-(\prod_{k=i}^m w_k)h_1\), proving the second. The scalar 1−W
is nonzero in the integral domain R, giving (C9). Formula (C10) follows
by applying J\_p and φ\_0 and then using both homotopies. For m=1 the
formulas give H=0 and J\_p=(1−w\_1); the two cyclic contributions are in
the same matrix entry. ∎

For completeness, the position-chain complex, rather than its dual, has
an equally explicit comparison with {[}R→R{]}. Send vertex 1 to 1,
vertex i≥2 to \(Q_i=\prod_{k=i}^m w_k\), edge 1 to 1, and other edges to
zero. The scalar sections send 1 to vertex 1 and to
\(\sum_iW_{i-1}E_i\). The homotopy sends vertex 1 to zero and vertex i≥2
to its weighted forward tail to vertex 1. Telescoping proves both
homotopy identities. These maps connect the two cyclic cokernel
descriptions through the same scalar complex; an unexplained equality of
a matrix and its transpose is not used.

\subsection{Proposition 3. The retained determinant and forcing
polynomial recover the word}

The pair (W,N\_p) determines the entire original labelled word. Its
inverse is explicit. Write \(W=u^m v^B\). In N\_p the coefficient of
u\^{}j is a single monomial \(v^{b_j}\), for 0≤j\textless m. Put
\(b_m=B\). The exact image conditions are
\SourceMath{b_0=0,\qquad 0=b_0\le b_1\le\cdots\le b_m,}{} with
coefficient 1 at each of these m original positions. The inverse is
\SourceMath{a_i=1+b_i-b_{i-1}\quad(1\le i\le m).}{C11}

\textbf{Proof.} By definition \(b_j=A_j-j\), giving every stated
condition and (C11). Conversely, (C11) yields positive integers with
cumulative exponents j+b\_j, whose total weight and forcing polynomial
are exactly the supplied pair. Both compositions are identities on the
specified image. ∎

The raw polynomial N\_p is retained, not only its quotient class. Merely
retaining W forgets precisely the original words with the same m,A. For
example, \SourceMath{W_{(1,3)}=W_{(2,2)}=u^2v^2,
\qquad N_{(1,3)}=1+u,\quad N_{(2,2)}=1+uv.}{C12} Thus the determinant
alone has a nontrivial word fiber, and (C11) supplies the actual repair.
This faithful marker uses the original weights of {[}W{]}; no priority
claim is made for polynomial encoding of finite words.

\section{3. The original integral cycle obstruction is an exact
specialization}

For the same word retain
\SourceMath{L=3^m,\quad U=2^{A_m},\quad D=U-L,\qquad
 C=\sum_{i=1}^m3^{m-i}2^{A_{i-1}},}{}
\SourceMath{(B_px)_i=2^{a_i}x_{i+1}-3x_i.}{C13} Let
\(\Lambda=\mathbb Z[1/3]\). The specified coefficient homomorphism is
\SourceMath{\alpha:R\longrightarrow\Lambda,\qquad u\longmapsto2/3,\quad v\longmapsto2.}{C14}
It sends \(w_i\) to \(2^{a_i}/3\). Directly, including all original
signs and factors, \SourceMath{\alpha(J_p)=-B_p/3,\qquad
 \alpha(1-W)=-D/3^m,\qquad
 \alpha(N_p)=C/3^{m-1}.}{C15} Thus the original forced system B\_px=b
becomes \(\alpha(J_p)x=-b/3\); its scalar forcing mark in (C6) becomes
\SourceMath{\alpha(\phi_1)(-b/3)=-F_p(b)/3^m,
 \quad F_p(b)=\sum_{i=1}^m3^{m-i}2^{A_{i-1}}b_i.}{C16} Multiplication by
the specified unit \(-3^m\) identifies the scalar equation with
\SourceMath{D t=F_p(b).}{C17} The degree-zero map from the localized
original complex to the evaluated position complex is the identity, and
its degree-one map is multiplication by −1/3. Its inverse is the
identity in degree zero and multiplication by −3 in degree one. This is
the chain isomorphism responsible for (C15)--(C17).

\subsection{Proposition 4. Localization at 3 retains the original
integral forcing class}

For original integral b, solvability of B\_px=b in \(\Lambda^m\) is
equivalent to solvability in \(\mathbb Z^m\). The map on original
cokernels is the explicit isomorphism
\SourceMath{\mathbb Z/D\mathbb Z\longrightarrow\Lambda/D\Lambda.}{C18}
It preserves the distinguished original class \([F_p(b)]\), in
particular {[}C{]} for b=1. It also retains the reduced denominator of
the unique rational primitive.

\textbf{Proof.} The integer D is odd and coprime to 3, since U is a
power of 2 and L a power of 3. The telescoping row F\_p satisfies
\(F_pB_p=D e_1^T\), is primitive (its endpoint coefficients are coprime
powers of 3 and 2), and both cokernels have determinant order
\textbar D\textbar{} before localization. Its explicit integral
identification is {[}C{]}'s calculation and {[}S{]}'s integral block
comparison. For (C18), the inverse sends \(a/3^r\) to
\(a(3^r)^{-1}\bmod D\). Independence of the chosen fraction and both
inverse identities follow by multiplication by powers of 3, invertible
modulo D.

An integral b with a localized preimage has 3\^{}r b in B\_p Z\^{}m for
some r. Its original cokernel class is killed by 3\^{}r, hence is zero
in Z/D by coprimality. An integral preimage therefore exists and equals
the localized one, since B\_p is injective. Equivalently its unique
rational coordinates are integer adjugate numerators divided by D, and
localization at 3 cannot remove a denominator coprime to 3. ∎

The coefficient embedding Z→Λ is not replaced by an identity on
lattices. Its extra degreewise quotient is the 3-primary group
\((\Lambda/\mathbb Z)^m\). On this group B\_p is an automorphism: its
adjugate multiplied by the inverse of D is an explicit inverse, because
D acts invertibly on 3-primary torsion. Hence that quotient complex is
acyclic. This accounts for the extra localized elements through their
actual quotient and maps.

The original positive cycle criterion remains D\textgreater0 and
D\textbar C; the full original coordinates and actual valuations are
reconstructed as in {[}C,S{]}. The ring change alone does not assert
positivity or the existence of an integer cycle. The word (1,3) has
(D,C)=(7,5), whereas (2,2) has (7,7). Their identical clock determinants
in (C12) give different original forcing classes through (C16). The
second word maps to the actual fixed vertex 1 and is removed only by the
declared relative graph quotient.

The prime-power control (1,1,2,4,3) retains D=1805 and C=475. Its scalar
equation has 19 solutions modulo 19 and none modulo 361, because
gcd(1805,19)=19 divides 475 while gcd(1805,361)=361 does not. Formula
(C18) and {[}S{]}'s integral inverse preserve that obstruction; no
passage to rational coefficients substitutes for this test.

\section{4. What the clocked cokernel records for actual components}

Let \(\mathscr C\) be the set of actual nonbase periodic components,
each equipped with one chosen point on its primitive cycle. Let
\(\mathscr A\) be the set of actual nonperiodic components, each
equipped with a chosen original forward ray \(x_0,x_1,\ldots\). These
are sets defined by the actual map, not assertions that either kind of
nonbase component exists.

For a primitive cycle c let m\_c,A\_c be its actual clocks and
\(W_c=u^{m_c}v^{A_c-m_c}\). For a chosen nonperiodic ray let
\SourceMath{B_j=\sum_{i<j}(a(x_i)-1),\qquad W_j=u^jv^{B_j},
\qquad M_x=\bigcup_{j\ge0}W_j^{-1}R
 \subset\mathbb Z[u^{\pm1},v^{\pm1}].}{C19} The union is increasing
because W\_j divides W\_\{j+1\}. Each element remains a finite Laurent
polynomial; the union is not silently replaced by a completion.

\subsection{Theorem 5. Component classification with explicit maps}

The cohomology of the actual relative complex is
\SourceMath{H^0(\mathcal K)=0,\qquad
 H^1(\mathcal K)\cong
 \bigoplus_{c\in\mathscr C} R/(1-W_c)
 \ \oplus\ \bigoplus_{x\in\mathscr A}M_x.}{C20} The original base
component contributes zero. Every cycle summand has its displayed
annihilator ideal (1−W\_c); every nonperiodic summand is a nonzero
torsion-free R-module inside the stated Laurent domain.

\textbf{Proof.} Finite-support modules decompose by actual components,
and d respects this decomposition. Proposition 1 gives H\^{}0=0.

On the component reaching 1, the actual finite path from n to 1 has
boundary V\_n by (C3), so every generator of the cokernel vanishes.
Equivalently, before taking the relative quotient, this component has
exactly R/(1−uv), and the subcomplex at 1 maps onto that summand by its
original generator. Taking the relative quotient removes precisely it.

For a periodic component, use the weighted finite path from every vertex
to the chosen cycle point to express its class as a monomial times that
point's class. One traversal of the primitive cycle gives exactly the
relation (1−W\_c)V=0. Send each vertex to the weight of its first path
to the chosen point, reduced modulo (1−W\_c). Every edge relation maps
to zero; at the chosen point it is precisely 1−W\_c.~The reverse map
sends 1 to the chosen vertex. The path relations and the primitive
return prove both inverse identities. This gives the displayed quotient,
whose annihilator is its defining principal ideal.

For a nonperiodic component, every vertex eventually meets the chosen
ray. To see this, adjacent vertices have forward orbits that meet, and
this property is transitive along each finite undirected path. For
T\^{}d(n)=x\_j, send
\SourceMath{[V_n]\longmapsto\frac{\prod_{i<d}w_{T^i(n)}}{W_j}\in M_x.}{C21}
Changing the meeting point to a later point multiplies numerator and
denominator by the same actual path weight. Nonperiodicity ensures that
two meeting descriptions align in this way, so the map is well defined.
It respects the edge relation by the first weight factor. Conversely
send f/W\_j to f{[}V\_\{x\_j\}{]}. Equality of two fractions can be
tested at a common later denominator W\_k, and the original ray path
relations make their images equal. These are inverse maps on generators,
and therefore on all finite sums. The target is a nonzero submodule of a
domain, proving torsion-freeness. ∎

This classification does not decide the actual component of a source
whose behavior has not been determined. It specifies exactly what a
component would contribute through the displayed maps, with no deletion
from a reference-measure statement.

\subsection{Repetitions and the actual orbit map}

For a solved positive integral word, the position vertex i maps to
V\_\{x\_i\} and edge i maps to E\_\{x\_i\}. The original equations prove
that the exponent weight is the actual w\_\{x\_i\}, so this is a chain
map for (C2). Its degreewise kernel consists exactly of coefficient
vectors whose sums over each repeated original vertex/edge label vanish.
The full coordinate dictionary is retained.

For r repetitions of a primitive cycle of weight W\_c, the source
cokernel maps by
\SourceMath{R/(1-W_c^r)\longrightarrow R/(1-W_c),\qquad[f]\longmapsto[f].}{C22}
Its kernel is \((1-W_c)/(1-W_c^r)\). The weighted fundamental path maps
to \SourceMath{(1+W_c+\cdots+W_c^{r-1})P_c.}{} At u=v=1 this gives
exactly r times the original cycle chain, not one traversal. Repetitions
of (2) map entirely to the removed base vertex and loop.

\subsection{Recovery of the old cycle kernel}

The coefficient homomorphism u→t,v→1 has kernel (v−1) and sends (C20)'s
cycle summand to Z{[}t{]}/(1−t\^{}\{m\_c\}), and its ray summand to
Z{[}t,t\^{}\{-1\}{]}. On a ray this follows directly by specializing the
explicit increasing union: the denominator W\_j becomes t\^{}j. The
kernel of the map on that union is (v−1)M\_x, as follows by writing an
element at one denominator and dividing its polynomial numerator by v−1.

Now use the exact sequence of complexes given by multiplication by t−1
and evaluation t=1. Its connecting map is explicit. For an ordinary
primitive cycle chain \(\sum_i E_i\), lift it to \(\sum_i t^{i-1}E_i\);
its differential is
\SourceMath{(1-t^m)V_1=-(t-1)(1+t+\cdots+t^{m-1})V_1.}{} Thus the
connecting class is \(-[1+t+\cdots+t^{m-1}]\), the generator of the
(t−1)-torsion subgroup in Z{[}t{]}/(1−t\^{}m). The Laurent ray module
has no such torsion. Quotienting either component module by t−1 gives Z.
This recovers exactly {[}C{]}'s ordinary cycle kernel and the
degree-zero component obstruction, with their signs and degree
correspondence accounted for.

\section{5. Two completion maps and the support they retain}

Let \(E_0=\mathbb Z[v]^{(X\setminus\{1\})}\) and
\(\widehat R=\mathbb Z[v][[u]]\). There is an explicit injective map
\SourceMath{\iota:\widehat R^{(X\setminus\{1\})}\longrightarrow E_0[[u]],
 \quad (f_n(u,v))_n\longmapsto\sum_{j\ge0}\left(\sum_n[u^j]f_n\,V_n\right)u^j.}{C23}
Its image is exactly the series with one common finite vertex-support
set for all u coefficients. The target permits the finite vertex support
to depend on the coefficient index. The quotient is precisely the latter
series modulo the common-finite-support series. These formulas specify
the relation between the two constructions, rather than merely giving
them different names.

On E\_0{[}{[}u{]}{]}, \SourceMath{(I-uS)^{-1}=\sum_{j\ge0}u^jS^j}{C24}
is defined coefficientwise. Each coefficient is a finite original vertex
sum. This is the forward-column form of {[}W{]}'s completed transfer
inverse, specialized to the relative source; {[}W{]} already contains
the existence of the completed inverse.

\subsection{Theorem 6. Exact source-support test for the inverse column}

For an original source n, the unique completed inverse column is
\SourceMath{G_n=\sum_{j\ge0}u^jv^{A_j(n)-j}V_{T^j(n)},\qquad V_1=0.}{C25}
The formula stops contributing after the first hit of 1. Its membership
has the following complete description:

\begin{itemize}
\tightlist
\item
  On an orbit reaching 1, G\_n is a polynomial vector, with u-degree
  exactly its positive number of odd returns to 1 minus one (and the
  zero column for n=1).
\item
  On an orbit eventually entering a nonbase cycle, G\_n lies in the
  image of (C23) and is not a polynomial in u. Its actual tail is a
  finite weighted cycle numerator divided by 1−W\_c, with its original
  entry prefix retained.
\item
  On a nonperiodic orbit, G\_n is outside the image of (C23): its
  nonzero coefficients have infinitely many distinct original vertex
  labels.
\end{itemize}

Consequently \SourceMath{[V_n]=0\text{ in }H^1(\mathcal K_{u,v})
 \quad\Longleftrightarrow\quad
 G_n\text{ is a polynomial vector}
 \quad\Longleftrightarrow\quad n\text{ reaches }1.}{C26}

\textbf{Proof.} Coefficient comparison in (I−uS)G\_n=V\_n uniquely gives
the displayed successive terms, retaining both clocks and the actual
label. Arrival at 1 makes all subsequent relative terms zero. A repeated
nonbase vertex makes the deterministic orbit periodic from its first
repetition; the finite prefix and geometric cycle tail give the stated
rational expression. Without any repeated vertex, the support is
unbounded in vertex labels. Conversely, finitely many encountered vertex
labels force a repetition or a base hit. These statements prove every
support characterization and (C26), since Proposition 1 makes a
polynomial preimage unique and its inclusion in the completion must
equal G\_n.~∎

The scalar extension from R to \(\widehat R\) does not by itself produce
(C24) on the common-finite-support module: on a ray its cokernel is the
nonzero increasing union \(\bigcup_j W_j^{-1}\widehat R\), with the same
fraction maps as (C21). For a finite cycle, 1−W\_c is a unit in
\(\widehat R\), so that cycle summand becomes zero. Passing further
through (C23) makes the ray column available and the entire completed
complex acyclic. This computes what each map removes.

Also, evaluation u=1 on polynomials does not extend to a unital
homomorphism \(\mathbb Z[v][[u]]\to\mathbb Z[v]\) fixing v: 1−u is a
unit in the source and would map to zero. The exact inclusion of
polynomials and this unit obstruction explain why an available formal
inverse is not automatically a finite-support contraction at the
original parameter value.

The global Collatz target is therefore the source-dependent polynomial
membership in (C26) for every n.~No uniform degree bound over all
positive integers is required or asserted. A completed inverse already
exists without that membership; its existence is not the missing proof.

\section{6. Retraction and root splicing retain the original clocks}

Let h,Q,F be one of {[}C,S{]}'s actual path retractions. Retain the
original path from n to its chosen root r(n), its length \(\ell_n\), and
its full exponent word. Let P\_n be its chain (C3) and M\_n its terminal
monomial. Define
\SourceMath{h_w(V_n)=P_n,\qquad Q_w(V_n)=M_nV_{r(n)},\qquad
 F_w=I-h_wd.}{C27} Root paths have length zero, and the relative base
vertex is zero. These are maps on the original finite-support polynomial
modules. The predecessor's positive edge multiplicities and the original
starting point reconstruct the path's chronological order: iterate its
deterministic successor for the sum of the multiplicities, then verify
both the endpoint and the complete edge-count dictionary. The code does
exactly that rather than inventing an ordering of a sparse vector.

\subsection{Theorem 7. Clocked homotopy and extension identities}

The maps satisfy
\SourceMath{dh_w=I-Q_w,\qquad h_wQ_w=0,\qquad Q_w^2=Q_w,}{}
\SourceMath{F_w^2=F_w,\qquad dF_w=Q_wd.}{C28} For an extension attached
at old roots as in {[}S{]}, let \(\Gamma_r\) be its full weighted path
from an old root to a final old root, including each old endpoint path.
Define \SourceMath{K_w(V_n)=M_n\Gamma_{r(n)},\qquad h'_w=h_w+K_w.}{C29}
Then \SourceMath{dK_w=Q_w-Q'_w,\quad F_wK_w=K_w,
 \quad Q'_wQ_w=Q_wQ'_w=Q'_w,
 \quad F'_wF_w=F_wF'_w=F'_w.}{C30} At u=v=1 these are exactly the old
retraction and root-splice maps.

\textbf{Proof.} Equation (C3) proves the first equation of (C28). The
maps h vanish at their roots, and their root monomials are 1, proving
the other vertex identities. Then \(h_wdh_w=h_w(I-Q_w)=h_w\), yielding
edge idempotence and the chain equation. In a concatenation, the second
path chain is multiplied by the first path's terminal monomial. Thus
(C29) is the actual added chain and telescopes to dK\_w=Q\_w−Q'\_w.
Final roots are old roots with no added path, so K\_wQ\_w=K\_w,
K\_wQ'\_w=0 and h\_wQ'\_w=0. In particular h\_wdK\_w=0 and hence
F\_wK\_w=K\_w. The remaining equations follow by multiplying I−h\_wd and
I−(h\_w+K\_w)d and applying those identities. The evaluation
homomorphism u=v=1 sends every prefix monomial to 1, recovering the
original chains coefficient for coefficient. ∎

For a genuine weighted cycle path P\_c based at n,
\SourceMath{F_wP_c=P_c-(1-W_c)h_w(V_n).}{C31} This is the exact relation
before specialization. At u=v=1, it becomes literal equality with the
old closed cycle chain. One must not assert that a generic weighted
cycle path has zero boundary: its boundary is exactly (1−W\_c)V\_n.

\subsection{An actual omitted-clock defect and its repair}

Use the old one-edge descent 5→1, with exponent 4. At the old root 7 the
new original block is
\SourceMath{7\xrightarrow{1}11\xrightarrow{1}17\xrightarrow{2}13\xrightarrow{3}5.}{}
It has terminal monomial \(u^4v^3\). The correct new column is
\SourceMath{P_{7\to5}+u^4v^3 E_5.}{} Its boundary is V\_7. Appending
E\_5 without that monomial produces instead
\SourceMath{V_7+(1-u^4v^3)V_5.}{C32} The verifier computes this residual
literally, then checks the corrected path and both compositions in
(C30). This is a concrete example in which retaining a clock is required
for the equation to hold.

\section{7. A reversible source database closes the 403-rise sector}

The previous retained sector was m=971,A=1539 with 403 exponent-1
positions and possible minimum in {[}99781,330749{]}. The new
certificate retains every positive odd source 3≤n≤330749, its exact
first-descent word, the full affine cylinder and target progression, and
its source parameter.

The database uses a finite ordered word dictionary and one dictionary ID
for each source row. The source for row i is n=3+2i. For its exact word
p, the original cell is \SourceMath{n=r+2Ut,\qquad T^m(n)=q+2Lt,}{} with
all of r,U,q,L,C,D and the first-descent interval retained. Its decoder
returns \SourceMath{n=3+2i,\quad p=\text{dictionary}[\text{id}_i],\quad
 t=(n-r)/(2U),\quad y=q+2Lt.}{C33} The inverse encoder looks up that
exact word and checks the same original n,t,y. This is an inverse on the
specified valid source records; it is not an assertion that a count or
hash recovers a word. Gzip is only a byte container: decompression is
verified to recover the exact canonical JSON.

The declared labels (1) and (2,2) remain present with empty
first-descent fibers and zero source counts. They are not removed from
the dictionary. Counts may be placed in
\(G(\mathbb Z)=\mathbb Z\sqcup\{\tau\}\) by assigning an integer to each
declared label and τ to an undeclared label. The reflection fixes each
integer and maps τ to 0. Its zero fiber is exactly \{τ,0\}; the separate
declaration mask distinguishes them. The full word, cell and original
source maps are retained in addition to this mask, so neither an
integral forcing class nor a vanished coefficient is inferred from the
support label alone.

\subsection{Theorem 8. Executed source-complete finite polynomial
contraction}

The delivered database and independent integer replay establish that
every positive odd n≤330749 reaches 1. They also construct a closed
actual witnessed graph with 262615 nonbase vertices, on which the
original unit edge cochain has an explicitly reconstructed nonnegative
integral primitive of maximum 164. On that witnessed graph, the clocked
operator satisfies \SourceMath{S^{164}=0,\qquad
 (I-uS)^{-1}=\sum_{j=0}^{163}u^jS^j.}{C34} Here
\(S V_n=v^{a(n)-1}V_{T(n)}\), with V\_1=0, so every coefficient and
original edge remains in (C34).

\textbf{Proof and executable certificate.}
\nolinkurl{sector_certificate.py} generates the actual first-descent
word for each of the 165374 source rows, retaining any failure rather
than accepting it. The saved database has 13989 occupied word cells and
the two present zero cells. \nolinkurl{independently_verify} replays the
original +1 operation using repeated integer division, separately from
the generator's valuation implementation. It checks every exponent,
absence of an earlier descent, the original affine source/target
parameter inverses, every complete cell's count on the interval, and the
coverage of every original source row. There are 575110 replayed
original return equations. Every target is a strictly smaller positive
odd integer, so strong induction beginning at 1 supplies convergence for
each seed.

More explicitly, assign height 0 to 1 and work through the source rows
in increasing order. For a first-descent path ending at y\textless n,
height(y) is already defined. Assign the intermediate vertices the
remaining path length plus height(y); all overlaps must agree and are
checked. The resulting original edge table is closed under the actual
map except for the explicitly included basepoint, and
\nolinkurl{height(n)−height(T(n))=1} is verified on all 262615 nonbase
vertices. Its maximum is 164; its largest original vertex is 8216025965.
Hence every sequence of 164 original edges reaches the relative zero
vertex. This proves the first identity in (C34) on every basis column.
Multiplication of the finite geometric sum proves both inverse
identities. The bound concerns this finite witnessed graph, not all
original positive integers. ∎

The new source interval includes 115485 additional odd seeds above the
prior bound 99779. It closes the entire old 403-rise sector without
enumerating its gap tuples: the possible minimum of each tuple is an
actual source in the certified interval, and therefore cannot lie on a
nonbase cycle.

\subsection{Theorem 9. All-length exclusion through 678 exponent-1
positions}

No nontrivial positive integer cycle has at most 678 exponent-1
positions in its primitive odd-return word.

\textbf{Proof.} Theorem 8 gives the minimum lower bound s=330751. For an
actual positive cycle,
\SourceMath{3^m<2^A=\prod_i(3+1/x_i)\le(3s+1)^m/s^m.}{} With k
exponent-1 positions, the unchanged identity
\(A=2m-k+\sum_{a_i\ge3}(a_i-2)\) implies
\SourceMath{(4s)^m\le2^k(3s+1)^m.}{C35} The certificate computes the
least power of 2 strictly above 3\^{}m by repeated doubling and verifies
\SourceMath{2^{\lceil\log_2(3^m)\rceil}s^m>(3s+1)^m
 \quad(1\le m\le1635).}{C36} The logarithm is notation for that exact
least-power integer; no floating-point logarithm is used. It also
verifies \SourceMath{(4s)^{1634}>2^{678}(3s+1)^{1634}.}{C37} Because
4s\textgreater3s+1, (C37) persists at every larger length and for all
k≤678. Thus (C35) forces m≤1633, while (C36) excludes every such length.
All finite comparisons and every source premise are regenerated by the
delivered checker. ∎

This is an internal, self-contained workbench bound, not a claim to
exceed published Collatz cycle bounds. In particular k counts exponent-1
positions, not the local-minimum count used in other cycle nomenclature.
The original clock map in {[}S{]} still gives local-minimum count m−k in
a nontrivial shortened cycle.

The next retained sector with k=679 is forced to
\SourceMath{m=1636,\quad A=2593,\quad
 330751\le\min x_i\le583287.}{C38} The checker verifies the two adjacent
minimum-endpoint inequalities, excludes the next power of 2 at that
period, and excludes larger lengths for k=679 by (C35). Then
\(\sum_{a_i\ge3}(a_i-2)=0\): the word contains 679 ones and 957 twos.
Cutting at the ones retains gap lengths \(r_i\ge0\), \(\sum r_i=957\),
the original rotation offset, and {[}S{]}'s anchored equations
\SourceMath{2\cdot4^{r_i}y_{i+1}-3^{r_i+1}y_i=2\cdot3^{r_i},\qquad y_i=x_i-1.}{C39}
The complete integral inverse in {[}S{]} and the marked clock
specialization in Section 3 both remain attached. Neither a solution nor
an exclusion of this entire next sector is asserted.

\section{8. What is new here, what is retained, and the outstanding
task}

The existing weighted sums, path category, ordinary graph homology,
affine packet identities and completed inverse are sourced, not
rediscovered under new names. The continuation adds the explicit marked
cyclic R-comparison and its original integral specialization; the
component cokernel and scalar-versus-support completion maps; the
clocked residual retraction and omitted-clock residual (C32); and the
reversible finite certificate establishing (C34)--(C37).

The original word is recoverable from (W,N\_p), the original integral
forcing is recoverable through (C14)--(C18), and actual component
behavior is attached through original integer vertices. The available
formal inverse is kept in its actual completion. This identifies the
productive role of retained data without a priority claim about an
undefined acronym or a claim that no other method can prove these
results.

The global remaining task is polynomial membership of G\_n for each
original positive source, or an actual nonzero original source class
with a certified orbit obstruction. Merely constructing G\_n in the
completion has already been done and does not settle that membership.
The next finite sector is (C38)--(C39); the old infinite nonperiodic and
first-crossing exception sources are unchanged. Zero reference mass
never removes one of those vertices.

This session used authenticated GitHub reads and the uploaded source
archive. The accessible GitHub tool surface exposed read operations
only; no write operation or CLI connection was available. Consequently
this continuation is supplied as an additive patch against the
authenticated head, not described as a pushed commit, PR update, or
successful new GitHub CI run. The original eight replays and the new
ordinary/optimized pair were executed locally; their receipts
distinguish that from the earlier remote CI.

\section{9. Verification controls}

The symbolic suite checks 364 full-word fixtures: every word of lengths
1 through 5 with letters 1 through 3, and the original five-letter
prime-power fixture. Each homotopy is checked as a polynomial matrix
identity on its basis columns; the arithmetic evaluation is checked
against the unchanged original B matrix. The source database is
independently replayed and every saved cell is compared with the
untouched predecessor constructor. The suite has 33503 named
algebraic/source checks in addition to its separately counted 575110
original return equations; the counts are not advertised as independent
mathematical reviews.

The periodic positive control is T\_−(n)=(3n−1)/2\^{}\{ν₂(3n−1)\}, with
the genuine orbit 5→7→5 and weight u²v. Its exact relation to the signed
original map is T\_−(n)=−T\_+(−n), using ν₂(−q)=ν₂(q). It is not a
positive +1 Collatz counterexample. Its r-fold weighted path is checked
against (1+u²v+⋯+(u²v)\^{}\{r−1\}) times the primitive path, for r≤16.

The ray control is expressly the synthetic map n→n+2, assigned the test
weight uv. Its length-j boundary is V\_3−u\^{}jv\^{}jV\_\{3+2j\}, with j
distinct original test labels. This checks the support-completion
mechanism; it is never assigned to an unverified +1 orbit. Twelve
malformed inputs are rejected, including an erased forcing coefficient,
a removed original source row, a changed forcing sign, and erased
supported-zero labels.
